% Arquivo TikZ refinado para reproduzir o framework original.
% Requer no preâmbulo do documento principal:
% \usepackage{tikz}
% \usetikzlibrary{arrows.meta,calc,positioning}

\begin{tikzpicture}[
    x=1cm,
    y=1cm,
    >=Stealth,
    line cap=round,
    line join=round,
    % font=\sffamily\small,
    arrow/.style={-{Stealth[length=3.2mm,width=3.8mm]}, line width=1.05pt},
    conn/.style={line width=1.25pt},
    panel/.style={draw=red!88!black, line width=3pt},
    model/.style={
        rectangle,
        draw=blue!45!black,
        fill=blue!18,
        minimum height=1.48cm,
        align=center,
        % font=\sffamily\bfseries\small,
        inner xsep=6pt,
        inner ysep=5pt
    },
    greenblock/.style={
        rectangle,
        draw=green!45!black,
        fill=green!18,
        align=center,
        % font=\sffamily\bfseries\small,
        inner xsep=5pt,
        inner ysep=4pt
    },
    title/.style={
        % font=\sffamily\bfseries\itshape\Large,
        align=center
    }
]

% ------------------------------------------------------------------
% Caixas com canto dobrado, no estilo do draw.io.
% Uso: \paperbox{nome}{(x,y)}{texto}{largura}{altura}
% ------------------------------------------------------------------
\newcommand{\paperbox}[5]{%
    \node[rectangle, draw=black, fill=white, align=center,
          minimum width=#4, minimum height=#5,
          inner xsep=4pt, inner ysep=4pt] (#1) at #2 {#3};
    \path (#1.north east) coordinate (#1ne);
    \draw[draw=black, fill=white]
        ($(#1ne)+(-0.34cm,0)$) -- (#1ne) --
        ($(#1ne)+(0,-0.34cm)$) -- cycle;
}

\newcommand{\greenpaperbox}[5]{%
    \node[greenblock, minimum width=#4, minimum height=#5] (#1) at #2 {#3};
    \path (#1.north east) coordinate (#1ne);
    \draw[draw=green!45!black, fill=green!18]
        ($(#1ne)+(-0.36cm,0)$) -- (#1ne) --
        ($(#1ne)+(0,-0.36cm)$) -- cycle;
}

% ------------------------------------------------------------------
% Molduras dos três módulos.
% ------------------------------------------------------------------
\draw[panel] (0.00,8.70) rectangle (7.75,14.25);
\node[title] at (3.90,13.88) {(i) Queue Management};

\draw[panel] (0.00,0.00) rectangle (14.65,7.45);
\node[title] at (7.30,7.06) {(ii) Overtime and Cancellation Management};

\draw[panel] (15.65,0.40) rectangle (24.75,7.30);
\node[title] at (20.20,6.90) {(iii) Operating Theatre Schedule};

% ------------------------------------------------------------------
% (i) Queue Management.
% Os blocos brancos foram espaçados verticalmente como no módulo (iii).
% ------------------------------------------------------------------
\paperbox{demand}{(1.20,12.65)}{Patient\\demand\\probability\\distribution}{1.50cm}{1.42cm}
\paperbox{waiting}{(1.20,10.90)}{Waiting time\\requirement}{1.50cm}{1.42cm}
\paperbox{confidence}{(1.20,9.15)}{Confidence\\level}{1.50cm}{1.42cm}

\node[model, minimum width=2.55cm] (rq) at (3.95,10.90) {$(R,Q)$ Control\\Policy};
\greenpaperbox{upper}{(6.55,11.78)}{Upper\\bound Q}{1.38cm}{1.20cm}
\greenpaperbox{lower}{(6.55,10.02)}{Lower\\bound Q}{1.38cm}{1.20cm}

% Entrada para a política (R,Q).
\draw[conn] (demand.east) -- (2.28,12.65) -- (2.28,10.90);
\draw[conn] (confidence.east) -- (2.28,9.15) -- (2.28,10.90);
\draw[arrow] (waiting.east) -- (rq.west);
\draw[arrow] (2.28,10.90) -- (rq.west);

% Saídas upper/lower bound Q.
\draw[conn] (rq.east) -- (5.38,10.90);
\draw[arrow] (5.38,10.90) |- (upper.west);
\draw[arrow] (5.38,10.90) |- (lower.west);

% Junção das duas saídas Q para seguir ao Preprocessing.
\draw[conn] (upper.east) -- (11.25,11.78) -- (11.25,10.90);
\draw[conn] (lower.east) -- (11.25,10.02) -- (11.25,10.90);

% ------------------------------------------------------------------
% (ii) Overtime and Cancellation Management.
% Blocos brancos com o mesmo ritmo vertical do módulo (iii).
% ------------------------------------------------------------------
\paperbox{surgdur}{(1.20,5.85)}{Surgery\\duration\\probability\\distribution}{1.50cm}{1.42cm}
\paperbox{otcost}{(1.20,4.15)}{Overtime\\cost}{1.50cm}{1.42cm}
\paperbox{idlecost}{(1.20,2.45)}{Idle time\\cost}{1.50cm}{1.42cm}
\paperbox{blocksize}{(1.20,0.75)}{Block size}{1.50cm}{1.42cm}

\node[model, minimum width=2.55cm] (newsvendor) at (4.25,3.30) {Newsvendor\\model};
\greenpaperbox{fitprob}{(7.45,4.15)}{Probabilities\\of surgeries\\to fit into the\\block}{1.95cm}{1.58cm}
\greenpaperbox{cancelprob}{(7.45,1.90)}{Probabilities\\of canceling\\the surgery}{1.95cm}{1.58cm}
\node[model, minimum width=2.65cm] (inventory) at (10.55,3.30) {Inventory-\\based model};
\greenpaperbox{numsurgeries}{(13.25,3.30)}{Number of\\surgeries to\\plan in a\\block}{1.65cm}{1.32cm}

% Entradas para o newsvendor.
\draw[conn] (surgdur.east) -- (2.55,5.85) -- (2.55,3.30);
\draw[conn] (otcost.east) -- (2.55,4.15);
\draw[conn] (idlecost.east) -- (2.55,2.45);
\draw[conn] (blocksize.east) -- (2.55,0.75) -- (2.55,3.30);
\draw[arrow] (2.55,3.30) -- (newsvendor.west);

% Newsvendor para probabilidades.
\draw[conn] (newsvendor.east) -- (5.95,3.30);
\draw[arrow] (5.95,3.30) |- (fitprob.west);
\draw[arrow] (5.95,3.30) |- (cancelprob.west);

% Probabilidades para modelo inventory-based.
\draw[conn] (fitprob.east) -- (8.85,4.15) -- (8.85,3.30);
\draw[conn] (cancelprob.east) -- (8.85,1.90) -- (8.85,3.30);
\draw[arrow] (8.85,3.30) -- (inventory.west);
\draw[arrow] (inventory.east) -- (numsurgeries.west);

% ------------------------------------------------------------------
% Preprocessing: junção limpa das duas fontes de entrada.
% ------------------------------------------------------------------
\node[model, minimum width=3.05cm, minimum height=0.78cm] (preproc) at (17.35,11.58) {Preprocessing};

% A saída do módulo (i) chega por uma linha horizontal superior.
\draw[conn] (11.25,10.90) -- (15.60,10.90) -- (15.60,11.78);

% A saída do módulo (ii) sobe e chega por uma linha horizontal inferior.
\draw[conn] (numsurgeries.east) -- (14.85,3.30) -- (14.85,11.38);

% Dois pequenos trechos finais com setas separadas, evitando cruzamento/ambiguidade.
\draw[arrow] (15.60,11.78) -- ([yshift=0.20cm]preproc.west);
\draw[arrow] (14.85,11.38) -- ([yshift=-0.20cm]preproc.west);

% ------------------------------------------------------------------
% (iii) Operating Theatre Schedule.
% ------------------------------------------------------------------
\paperbox{compat}{(16.95,4.85)}{Theatre\\compatibility}{1.68cm}{1.42cm}
\paperbox{teams}{(16.95,3.15)}{Teams\\availability}{1.68cm}{1.42cm}
\paperbox{anaesth}{(16.95,1.45)}{Anaesthetists\\availability}{1.68cm}{1.42cm}

\node[model, minimum width=2.55cm] (milp) at (20.70,3.15) {MILP allocation\\model};
\greenpaperbox{monthly}{(23.35,3.15)}{Monthly\\Theatre\\Schedule}{1.35cm}{1.18cm}

% Preprocessing para o MILP.
\draw[conn] (preproc.south) -- (17.35,6.05) -- (18.62,6.05) -- (18.62,3.76);
\draw[arrow] (18.62,3.76) -- (milp.west);

% Restrições/entradas locais para o MILP.
\draw[conn] (compat.east) -- (18.35,4.85) -- (18.35,3.15);
\draw[conn] (anaesth.east) -- (18.35,1.45) -- (18.35,3.15);
\draw[arrow] (teams.east) -- (milp.west);
\draw[arrow] (18.35,3.15) -- (milp.west);

% Saída final.
\draw[arrow] (milp.east) -- (monthly.west);

\end{tikzpicture}
